\chapter{Interactive GSLTs and Behavioral Primes}
\label{sec:combinators}
% ===================================================================

\section{Interactive GSLTs}

\begin{definition}[Interactive GSLT]
A GSLT $\GSLT$ is \emph{interactive} if it has a designated
\emph{interaction operator} $\inter$ such that:
\begin{enumerate}[label=(\roman*)]
  \item $\inter : \terms(\GSLT) \times \terms(\GSLT) \to \terms(\GSLT)$
        is a term constructor (part of the grammar);
  \item the \emph{base rewrite rule} has the form
        $P \inter Q \rewrite R$ for terms $P, Q, R$,
        requiring no contextual hypothesis to fire;
  \item all other rewrite rules are \emph{context rules}: they carry
        hypotheses (in $\HML(\ctx)$) specifying which contexts
        permit the base rule to propagate.
\end{enumerate}
\end{definition}

\begin{example}
In the $\lambda$-calculus, $\inter$ is application: $P \inter Q = PQ$.
The base rule is $\beta$-reduction $(\lambda x.M)N \rewrite M\{N/x\}$.
The context rules determine reduction strategy
(call-by-value, call-by-name, etc.).
Application is \emph{asymmetric}: the left position is the program
(function), the right is the environment (argument).

In the Rho calculus, $\inter$ is parallel composition: $P \inter Q = P \mid Q$.
The base rule is the comm rule.
The context rules (par, struct) specify when and where comm fires.
Parallel composition is \emph{symmetric}: either term may be program
or environment depending on which carries the active receptor.
\end{example}

\begin{remark}[Program and Environment]
The asymmetry of $\inter$ is a real ontological distinction.
When $\inter$ is asymmetric, there is a canonical program/environment
boundary baked into the grammar.
When $\inter$ is symmetric, the boundary is emergent: it is drawn by
the rewrite rules at runtime, not by the grammar.
This distinction will matter for how agents perceive their world:
in a symmetric GSLT, any agent may find itself in either the program
or the environment role depending on context.
\end{remark}

\section{Behavioral Primes}

\begin{definition}[Combinator Set]
A \emph{combinator set} for an interactive GSLT $\GSLT$ is a set
$\comb \subseteq \terms(\GSLT)$ such that every term
$P \in \terms(\GSLT)$ is bisimilar to a finite interaction of
elements of $\comb$:
\[
  \forall P \in \terms(\GSLT),\;
  \exists c_1, \ldots, c_n \in \comb \text{ such that }
  P \bisim c_1 \inter c_2 \inter \cdots \inter c_n.
\]
\end{definition}

\begin{definition}[Behavioral Prime]
An element $c \in \comb$ is a \emph{behavioral prime} if it is not
bisimilar to any interaction $c' \inter c''$ of smaller terms.
A combinator set $\comb$ is \emph{prime} if every element is a
behavioral prime and no proper subset of $\comb$ is still a
combinator set.
\end{definition}

\begin{remark}[Computational Quarks]
The elements of a prime combinator set $\comb$ are the
\emph{computational quarks} of $\GSLT$: irreducible behavioral units
whose interactions generate all observable behavior.
They are \emph{not} the bisimulation equivalence classes, which may
be equivalence classes of arbitrarily complex composite behavior.
The primes are indivisible in the interaction-factorization sense:
they cannot be split into independently interacting parts.
\end{remark}

\begin{conjecture}[Existence of Prime Combinator Sets]
Every interactive GSLT with a finitely generated grammar admits a
prime combinator set.
\end{conjecture}

\begin{remark}
Unique factorization --- the analogue of the fundamental theorem of
arithmetic for behavioral primes --- is a stronger condition that
need not hold in general.
Its failure in a given GSLT would be physically meaningful:
it would correspond to composite objects that cannot be assigned a
unique quark content, analogous to the situation in QCD where
the quark content of hadrons is not always uniquely determined.
Whether unique factorization holds is a structural question about
the specific GSLT.
\end{remark}

% ===================================================================
